The fourth example demonstrates the engagement-classification stage that runs
after every assistant response containing a sponsored ad.  A user who has just
received the Nike Pegasus~41 recommendation follows up with
\emph{``How does it compare to the Asics Nimbus?''}---a question that engages
directly with the advertised product.  Figure~\ref{fig:example-4} shows the
chat interface at this point in the conversation.

\begin{figure}[H]
    \centering
    \includegraphics[width=\textwidth]{images/illustrative_example_4.png}
    \caption{Chat interface after the user asks a comparison question following
    the Nike Pegasus~41 sponsored recommendation.  The follow-up triggers the
    engagement classifier and a new engagement event is written to the database.}
    \label{fig:example-4}
\end{figure}

\subsubsection{Stage 1: Engagement Classification}
\label{subsubsec:example4-engagement}

Once the response is streamed, the orchestrator dispatches the prior assistant
turn and the new user message to the LLM-based engagement judge.  The judge
evaluates whether the follow-up message constitutes genuine engagement with the
sponsored product and, if so, characterises the nature of that engagement.  For
this turn it returns:

\begin{lstlisting}[style=jsonstyle, caption={Engagement classifier output for
the comparison follow-up.}, label={lst:example4-engagement}]
{
  "engaged":               true,
  "engagement_type":       "comparison",
  "engagement_score":      0.82,
  "ad_naturalness_score":  0.78,
  "purchase_proximity":    0.55,
  "sentiment_toward_ad":   "positive",
  "follow_up_topic_match": true,
  "reasoning": "User directly compares the advertised shoe to a competitor,
                indicating active product consideration."
}
\end{lstlisting}

\noindent
The \texttt{engaged} flag is set to \texttt{true} because the user stayed on
the product topic (\texttt{follow\_up\_topic\_match}: \texttt{true}) and asked
a question that can only arise from genuine interest in the Nike Pegasus~41.
The \texttt{engagement\_type} of \texttt{comparison} distinguishes this event
from a simple \texttt{product\_inquiry}: the user is actively evaluating
alternatives, which the classifier maps to a higher \texttt{purchase\_proximity}
of~0.55 than a general enquiry would warrant.  The \texttt{ad\_naturalness\_score}
of~0.78 reflects that the original recommendation felt sufficiently organic to
invite a follow-up rather than prompt dismissal.

\subsubsection{Stage 2: Event Persistence}
\label{subsubsec:example4-persistence}

The full classifier payload is written to the \texttt{events} table as an
\texttt{engagement} event with \texttt{event\_type = "engagement"} and the JSON
object stored in \texttt{events.payload}.  This record is the raw material for
all downstream analytics: the engagement-analytics endpoint aggregates these
payloads to compute per-ad averages of \texttt{engagement\_score},
\texttt{ad\_naturalness\_score}, and \texttt{purchase\_proximity}, and counts
occurrences of each \texttt{engagement\_type} to drive the engagement-type
breakdown chart.

\subsubsection{Stage 3: Analytics Reflection}
\label{subsubsec:example4-analytics}

After this event is logged the engagement-analytics page for the Nike
Pegasus~41 ad updates to reflect the new data point.  The three sub-metric
cards---engagement score, naturalness, and purchase proximity---each recalculate
their weighted averages.  The engagement-type pie chart gains or enlarges a
\emph{comparison} slice.  Because \texttt{purchase\_proximity} is now~0.55 (above
the 0.4 caution threshold), no intervention flag is raised; if it had fallen
below that threshold, the analytics view would highlight the ad as a candidate
for context-level optimisation via the mechanism described in
Example~\ref{subsec:example-optimization}.
